\documentclass[review]{elsarticle}
\usepackage[T1]{fontenc}
\usepackage{graphicx}
\usepackage{booktabs}
\usepackage{amsmath}
\usepackage{array}
\usepackage{longtable}
\usepackage{tabularx}
\usepackage{enumitem}
\usepackage[hidelinks]{hyperref}
\usepackage{url}
\graphicspath{{../figures/}}
\newcolumntype{Y}{>{\raggedright\arraybackslash}X}
\journal{Computers and Electronics in Agriculture}

\begin{document}
\begin{frontmatter}
\title{Multi-View and Multimodal Perception for Class-Wise Fruit Inventories: A Design-Space Review}
\author[a]{Muhammad Zainal Muttaqin\corref{cor1}}
\ead{mz.muttaqin1@gmail.com}
\author[a]{Fatma Indriani}
\affiliation[a]{organization={Department of Computer Science, Universitas Lambung Mangkurat},city={Banjarbaru},country={Indonesia}}
\cortext[cor1]{Corresponding author}
\begin{abstract}
Converting observations of a tree into one inventory of unique fruits grouped by class requires more than per-image detection. A system must identify instances, attach an attribute such as maturity or quality to the correct instance, associate repeated observations, and report the counting error at tree level. This design-space review addresses that task in agricultural fruit perception, using oil-palm fresh fruit bunches as the principal case while transferring mechanisms from non-agricultural work on tracking, re-identification, multi-view geometry, and 3D association. Scopus and OpenAlex searches were executed with seven query families covering fruit counting, cross-view identity, plant geometry, per-instance attributes, previous reviews, oil palm, and single-view fruit-load estimation. The exported records were deduplicated by DOI or normalized title and year, and high-confidence title and abstract exclusions were recorded before full-text retrieval. The search produced 21,066 master records; 20,035 remained for full-text assessment, but the local repository contained verified full text for only six matched records. The synthesis therefore separates search candidates from the verified evidence register. The resulting taxonomy has six identity mechanisms, from intra-view counting and statistical correction to appearance, temporal, geometric, and learned multi-view association. Existing reviews primarily organize ripeness methods, detectors, or sensors; the accessible review sample does not make cross-view identity the central unit of analysis. The main contribution is consequently a testable design question: under which acquisition assumptions does explicit deduplication improve class-wise inventory enough to justify its data, calibration, and computation requirements?
\end{abstract}
\begin{keyword}
multi-view perception \sep multimodal fusion \sep fruit counting \sep instance identity \sep oil palm \sep design-space review
\end{keyword}
\end{frontmatter}

% Final revision candidate based on the eight lecturer revision points.
% Search candidates and verified full-text evidence are kept as separate registers.

\section{Introduction}
\label{sec:intro2}

An agricultural vision system may observe the same tree from several positions, frames, or sensors. The operational output is not the number of detections emitted by those images. It is an inventory of unique fruit instances, grouped by an attribute class. For a tree with instance set $I$, the desired count for class $c$ is the number of instances whose assigned attribute is $c$. A detection that is missed, merged with a neighbour, assigned the wrong class, or counted again from another view produces a different error. These errors cannot be diagnosed from image-level accuracy alone.

This review addresses the task as a design problem:
\begin{quote}
\emph{How can many observations of one tree be converted into a class-wise inventory in which each fruit or fruit bunch is counted once?}
\end{quote}
The principal agricultural case is the oil-palm fresh fruit bunch (FFB). The class attribute is not restricted to maturity. Maturity is one example; size, quality, disease, defect, cultivar, or harvest readiness are also possible attributes. All of them remain instance-dependent because a class-wise count is undefined until an individual fruit or bunch has been localized and assigned an identity.

The scope is agricultural fruit perception, but the mechanism search is deliberately wider. Tracking, person or vehicle re-identification, multi-camera association, visual odometry, structure from motion (SfM), and 3D object detection were developed largely outside orchards. They are included as transferable evidence when they answer a mechanism question about identity, geometry, or duplicate rejection. Their results are not presented as agricultural performance.

The review makes four contributions. First, it reports a reproducible Scopus and OpenAlex search with query strings, date bounds, selection rules, and a deterministic deduplication record. Second, it organizes the task around six identity mechanisms and the acquisition assumptions each mechanism requires. Third, it positions previous oil-palm, fruit-counting, and fruit-ripeness reviews against the cross-view identity problem. Fourth, it provides a row-level synthesis matrix in which method, modality, evaluation context, identity mechanism, class attribute, finding, and transfer boundary are visible together.

The paper does not propose one architecture as the universal solution. It asks when a simple statistical correction, temporal tracking, appearance matching, geometric reconstruction, or learned multi-view model is justified by the acquisition regime. This distinction is important for oil palm because the same tree can be observed from four or eight positions while sensor calibration, depth quality, illumination, and occlusion vary at the same time.

\section{Review questions and reproducible protocol}
\label{sec:protocol2}

\subsection{Review type and questions}

This work is a structured design-space review reported with a PRISMA-style search and selection trace. It is not a meta-analysis. The included tasks, object definitions, sensor modalities, and evaluation metrics are too heterogeneous for a pooled effect size. The protocol was not registered before execution. The review questions are:

\begin{enumerate}[leftmargin=*,itemsep=1pt]
  \item What operations are required to turn repeated observations into a unique, class-wise inventory?
  \item Which identity mechanisms have been used for fruit or for transferable non-agricultural objects?
  \item Which assumptions about overlap, motion, calibration, appearance, and view order make each mechanism plausible?
  \item What evidence exists for oil-palm and other tree-fruit systems, and what remains a transfer hypothesis?
  \item Which measurements would distinguish a better detector from a better inventory system?
\end{enumerate}

\subsection{Sources, date range, and query families}

Scopus was used as the subscription index available in the captured institutional session. OpenAlex was used as an open complementary index. Web of Science was not available in the recorded run; this deviation is stated rather than hidden. Google Scholar was reserved for possible backward or forward snowballing and was not used as a primary count source because its result export is not stable. The date window was 2015--2026. In Scopus the upper and lower bounds were written explicitly as \texttt{PUBYEAR > 2014 AND PUBYEAR < 2027}; the equivalent Web of Science syntax is recorded in the protocol for a future rerun.

Seven query families were used. Q1 targets fruit inventories from multiple observations. Q2 targets cross-view identity and duplicate rejection, including non-agricultural objects. Q3 targets plant geometry and multimodal sensing. Q4 targets class attributes attached to instances. Q5 targets previous reviews for positioning. Q6 targets oil-palm vision and machine learning. Q7 targets single-view fruit counting and yield estimation, which Q1 would otherwise miss. The exact strings used in the Scopus run are reproduced in Appendix~\ref{app:queries2}. The canonical query file is \path{literature/search/scopus-queries.md}; its count export and SHA-256 fingerprints are in \path{literature/search/scopus-counts-2026-08-09.csv}.

\begin{table}[t]
\centering
\caption{Search yields before cross-source deduplication.}
\label{tab:searchcounts2}
\small
\begin{tabular}{lrrrr}
\toprule
Query & Scopus & OpenAlex & Main focus & Known item\\
\midrule
Q1 & 2,001 & 1,849 & multi-view fruit inventory & SawitMVC\\
Q2 & 2,789 & 1,991 & identity and duplicates & mechanism set\\
Q3 & 5,648 & 6,422 & plant geometry and sensors & Gené-Molá\\
Q4 & 2,422 & 2,757 & per-instance attributes & class labels\\
Q5 & 662 & 1,143 & previous reviews & review set\\
Q6 & 995 & 1,177 & oil-palm vision & FFB studies\\
Q7 & 1,205 & 1,317 & single-view fruit load & MangoYOLO\\
\midrule
Total exported & 15,722 & 16,656 & & \\
\bottomrule
\end{tabular}
\end{table}

The Scopus known-item checks passed for SawitMVC in Q1, Gené-Molá et al. in Q3, and MangoYOLO in Q7. The checks demonstrate that the subject-area filter did not silently remove the three anchors. They do not demonstrate that either index is complete.

\subsection{Eligibility and elimination rules}

The following rules were applied to title, abstract, and full text in sequence. A source was eligible for the focused evidence register when all five inclusion conditions were met:

\begin{description}[leftmargin=*,style=nextline,itemsep=2pt]
  \item[IC1.] It is a peer-reviewed research article, conference paper, or openly available preprint with a verifiable full text and a quantitative result.
  \item[IC2.] Its visual output is per instance, such as object detection, instance segmentation, tracking, localization, or cross-observation association.
  \item[IC3.] Its publication year is 2015--2026, except for foundational Register B citations used only for context.
  \item[IC4.] It is written in English or Indonesian.
  \item[IC5.] The full text identifies a dataset or acquisition protocol and an evaluation metric or quantitative comparison.
\end{description}

The elimination rules were EC1, output only at image or scene level without per-instance identity; EC2, no quantitative evaluation; EC3, full text unavailable; EC4, duplicate or version of an already retained record; EC5, editorial, abstract-only, poster, patent, or non-research document; and EC6, outside the date range. A title or abstract was excluded automatically only when the signal was high confidence. Ambiguous records were retained for full-text assessment.

\subsection{Deduplication and screening trace}

The raw files were kept unchanged. Within each query, exact duplicate identifiers were removed. Across sources, the primary key was a normalized DOI. Records without a DOI used normalized title plus year. A title-year collision was not merged when author or venue evidence conflicted. The resolver then merged only groups with strong author overlap or an explicit DOI and venue agreement; conservative groups remained separate when evidence was insufficient.

The seven Scopus exports contained 15,722 rows and the seven OpenAlex exports contained 16,656 rows. Together they contained 32,378 raw rows, 32,222 rows after within-query deduplication, and 22,269 records after exact-key union. The resolver processed 1,030 merge components and reduced the working master to 21,066 records. The complete decision records are stored in the following files:
\begin{quote}
\path{literature/search/derived/dedup-resolution-2026-08-09.csv}\par
\path{literature/search/derived/master-screening-2026-08-09.csv}
\end{quote}

High-confidence title screening removed 242 EC5 records and one EC6 record, leaving 20,823 records for abstract screening. Abstract screening removed 640 EC1 records and 148 EC5 records, leaving 20,035 for full-text assessment. Of these, 8,624 carried a strong mechanism signal and 11,411 were retained because the abstract was ambiguous or unavailable. The local PDF audit found 182 verified PDFs from the earlier evidence corpus, but only six matched the new master. All six were text-extractable; five had evidence of a per-instance or cross-observation mechanism, quantitative evaluation, and dataset or protocol, while one review record requires manual IC1, EC2, and EC5 confirmation. The remaining 20,029 full-text candidates are not represented as included studies in this manuscript.

\begin{figure}[t]
\centering
\small
\begin{tabular}{c}
\fbox{\parbox{0.78\columnwidth}{\centering 32,378 raw records\\Scopus 15,722 + OpenAlex 16,656}}\\[2pt]
$\downarrow$\\[-1pt]
\fbox{\parbox{0.78\columnwidth}{\centering 22,269 exact-key master records}}\\[2pt]
$\downarrow$\\[-1pt]
\fbox{\parbox{0.78\columnwidth}{\centering 21,066 resolved working-master records}}\\[2pt]
$\downarrow$\\[-1pt]
\fbox{\parbox{0.78\columnwidth}{\centering 20,035 full-text candidates\\6 matched local PDFs; 20,029 still unavailable locally}}
\end{tabular}
\caption{Search and screening trace at the reporting cutoff. The final box is a retrieval boundary, not a completed included-study count.}
\label{fig:flow2}
\end{figure}

\subsection{Two evidence registers}

The manuscript keeps two registers separate. Register A is the new Scopus and OpenAlex search, including all candidate and screening states. Register B is the focused synthesis register: studies whose full text or authoritative article record was available for a mechanism-level statement. This separation prevents the old 182-PDF corpus from being presented as if it were the endpoint of the new 21,066-record search. It also prevents an automated abstract triage from being described as a completed full-text decision.

\section{The inventory task as a design space}
\label{sec:designspace2}

\subsection{From detections to inventory}

The task has five distinct operations. First, a detector or segmenter proposes an instance in each observation. Second, a class head attaches an attribute to that instance. Third, observations are associated when they refer to the same physical object. Fourth, duplicate observations are rejected or merged. Fifth, the remaining identities are aggregated into counts per class. The first two operations can be evaluated with precision, recall, average precision, and class confusion. The last three require identity and count metrics.

This decomposition changes the interpretation of a strong detector. A detector with high average precision on independent images may still overcount when a tree is filmed from both sides. A tracker may preserve identity through a short video but fail when the camera returns after a gap. A 3D association method may remove duplicates but become unusable when depth is uncalibrated. The design choice is therefore conditional on the acquisition regime.

\subsection{Six identity mechanisms}

Table~\ref{tab:mechanisms2} defines the six mechanisms used in the synthesis. M0 is the no-association baseline. M1 uses a statistical correction or a fixed counting protocol without assigning a persistent identity. M2 matches local appearance or a learned re-identification embedding. M3 uses temporal continuity, motion, and track management. M4 uses calibrated geometry, depth, SfM, or a 3D map. M5 learns a multi-view association function directly from linked observations. A system can combine mechanisms, but its reported result must identify which mechanism is responsible for duplicate rejection.

\begin{table}[t]
\centering
\caption{Identity mechanisms and their minimum assumptions.}
\label{tab:mechanisms2}
\small
\begin{tabularx}{\columnwidth}{@{}lY Y@{}}
\toprule
Code & Mechanism & Main assumption and measurable risk\\
\midrule
M0 & Intra-view detection or segmentation & One view is sufficient. Hidden objects and repeated views are not resolved.\\
M1 & Statistical correction or fixed census & Duplicate structure is stable enough for a correction factor. The correction may not transfer across trees.\\
M2 & Appearance or re-identification & The instance retains discriminative texture, colour, or shape despite illumination and occlusion.\\
M3 & Temporal tracking & Consecutive observations overlap in time and motion is sufficiently smooth. Track fragmentation and drift remain possible.\\
M4 & Geometric or 3D association & Calibration, scale, or reliable reconstruction is available. Sensor failure and SfM cost are risks.\\
M5 & Learned multi-view association & Linked training data represent view, class, and occlusion variation. The model may overfit one acquisition pattern.\\
\bottomrule
\end{tabularx}
\end{table}

\subsection{Acquisition assumptions}

The second design axis is the assumption set. A1 is the amount of overlap and occlusion. A2 is the camera baseline or view separation. A3 is motion smoothness or scene stationarity. A4 is calibration and metric scale. A5 is appearance distinctiveness. A6 is regularity of duplicate exposure, for example a fixed four-side census. A7 is the availability of identity supervision or an oracle association. M1 mainly depends on A6. M2 depends on A5. M3 depends on A3. M4 depends on A2 and A4. M5 depends on A7 and on coverage of the training acquisition pattern.

This two-axis view gives a precise meaning to the phrase ``multi-view deduplication.'' It is not a mandatory module. It is a mechanism whose value should be tested against the assumptions that survive in the field. A fixed four-side protocol may make M1 competitive when duplication is regular and metric depth is unreliable. An unstructured moving camera may require M3 or M4 because the number and order of observations are not fixed.

\subsection{Class attributes beyond maturity}

Maturity classification is common in FFB literature, but the identity dependency is more general. An instance may carry a maturity label, a size class, a defect label, a disease state, a cultivar label, or a harvest-readiness state. The inventory should first establish one identity record and then attach one or more attributes to that record. If the same bunch appears in three views with two maturity predictions, the system must report whether the conflict is resolved by confidence, temporal evidence, geometry, or manual adjudication. Summing per-view class predictions without this step is not a class-wise inventory.

\section{Direct agricultural and fruit evidence}
\label{sec:direct2}

\subsection{What the oil-palm literature measures}

The direct oil-palm literature is concentrated on detection, ripeness classification, dataset construction, and yield-related counting. Lai et al. reviewed FFB ripeness detection methods across sensing families, while Goh et al. reviewed non-destructive classification methods including computer vision, LiDAR, spectroscopy, thermal sensing, and laser-light backscatter \cite{lai2023palmslr,goh2025palmreview}. These reviews establish that sensor choice and class definition are central problems. Their primary unit is ripeness assessment, not the persistent identity of one bunch across repeated views.

Naftali et al. compared YOLO and other detectors on an oil-palm plantation dataset and reported a depthwise-separable YOLOv8 variant as a real-time candidate \cite{naftali2024palmcounter}. This is useful evidence for M0 detector selection and deployment constraints. It does not by itself measure whether a bunch detected in one frame is the same bunch detected in another frame. A detector benchmark and an inventory benchmark answer different questions.

Three dataset papers expose the acquisition variables needed for the new design space. Goh et al. released 400 RGB images and 400 registered depth maps from four plantation locations, captured with an Intel RealSense D455f. The article reports a mean registration error of 1.8 cm at 3 m and 92.5\% expert inter-rater agreement for binary ripeness labels \cite{goh2025palmdataset}. Suharjito et al. released plantation imagery captured from multiple angles with five categories: unripe, underripe, ripe, flower, and abnormal; the dataset contains 10,207 training images, 2,896 validation images, and 1,400 test images after the reported augmentation and split procedure \cite{suharjito2025palmdataset}. These datasets address sensor and class variation, but their public descriptions do not make the same cross-view identity graph the primary evaluation object.

SawitMVC is directly aligned with the inventory problem. Its source article describes 3,992 smartphone images of 953 oil-palm trees from two plantations, captured from four or eight angular positions. The dataset uses four BBC maturity classes and a separate per-tree JSON layer with 9,823 unique bunches and cross-view identity links \cite{indriani2026sawitmvc}. This makes a tree-level unique-count metric possible. It also exposes the central confounder: the number of visible bounding boxes is not the number of unique bunches.

\subsection{Fruit and orchard evidence}

Single-view fruit-load studies show the limit of a detector-only interpretation. Koirala et al. compared six existing architectures and developed MangoYOLO for mango detection. Their best reported model achieved an F1 score of 0.968 and average precision of 0.983 on an independent test set, while orchard load estimates still required a correction factor against hand-harvest counts \cite{koirala2019deepfruit}. The correction is an M1 mechanism: it can be operationally useful without solving identity, but its validity depends on stable exposure and sampling assumptions.

Gené-Molá et al. made identity and 3D location explicit by combining Mask R-CNN instance segmentation with SfM, projecting image detections into a 3D point cloud, and removing false positives with a support vector machine. The study used 11 Fuji apple trees containing 1,455 apples. The reported F1 score increased from 0.816 for 2D detection to 0.881 for 3D detection and location, while the authors identified SfM processing time as a barrier to real-time use \cite{genemola2020fruit3d}. This is direct evidence that geometric association can change the counting result, but it is not evidence that SfM is the best field solution for oil palm.

Fu et al. used RGB and depth features with Faster R-CNN for apple detection in dense foliage and robotic harvesting \cite{fu2020apple}. The study supports the role of depth as an occlusion and localization cue. It does not settle whether depth remains reliable under bright outdoor oil-palm conditions. The new search also returned citrus and grape yield-estimation studies, including multi-period citrus estimation and a YOLO-based grape framework \cite{zhai2025citrusyields,ma2026grapeyield}. These records widen the agricultural comparison, but their exact full-text inclusion status remains tied to the retrieval boundary in Section~\ref{sec:protocol2}.

\subsection{Positioning against previous reviews}

The reviews found by Q5 are adjacent rather than identical to this review. Table~\ref{tab:position2} records the scope visible in the accessible article record and identifies whether cross-view identity is treated as the primary unit. The wording is deliberately bounded: ``not explicit'' means that the accessible scope and abstract do not make identity and duplicate rejection the central organizing question. It does not claim that no cited paper inside the review discusses tracking.

\begin{table}[t]
\centering
\caption{Positioning against previous review and survey records.}
\label{tab:position2}
\scriptsize
\begin{tabularx}{\columnwidth}{@{}p{0.25\columnwidth}p{0.21\columnwidth}p{0.18\columnwidth}Y@{}}
\toprule
Review & Primary scope & Cross-view identity central? & Position relative to this review\\
\midrule
Lai et al. \cite{lai2023palmslr} & Oil-palm FFB ripeness detection methods & Not explicit & Supplies direct agricultural sensing context; does not organize unique counting mechanisms.\\
Goh et al. \cite{goh2025palmreview} & FFB ripeness classification across non-destructive sensors & Not explicit & Supplies modality and deployment limitations; this review adds identity and inventory evaluation.\\
Naftali et al. \cite{naftali2024palmcounter} & Oil-palm detector and counter comparison & Not explicit & Supplies detector and edge-computation context; this review separates detector score from duplicate-free count.\\
Naghipour et al. \cite{naghipour2024fruitreview} & AI methods for fruit detection and classification & Not explicit & Supplies broad fruit and model coverage; this review makes cross-view association a design axis.\\
Sapkota et al. \cite{sapkota2024yoloagri} & YOLO variants in agricultural applications & Not explicit & Supplies architecture lineage and application taxonomy; this review is task-first rather than model-first.\\
Nordin and Misro \cite{nordin2026palmslr} & Deep-learning review of palm-oil maturity detection & Not explicit in accessible record & Adds a recent maturity-focused comparator; this review generalizes the class attribute beyond maturity.\\
\bottomrule
\end{tabularx}
\end{table}

The bounded conclusion from Q5 is therefore not ``no previous review exists.'' It is that the accessible review set is organized mainly by sensor, ripeness, detector family, or agricultural application. The cross-view identity of each physical fruit as the prerequisite for a class-wise inventory is not the primary comparison unit in that set. This is the gap addressed here, subject to the stated search and retrieval boundary.

\section{Transferable mechanisms for identity and duplicate rejection}
\label{sec:transfer2}

\subsection{Appearance and temporal association}

M2 matches instances by appearance. A descriptor may use colour, texture, shape, or a learned embedding. Its strength is that it does not require metric depth. Its failure modes are predictable in fruit scenes: two bunches of the same maturity can look alike, illumination changes colour, and partial occlusion removes the most distinctive region. An evaluation must therefore report identity precision and recall under appearance changes, not only detector average precision.

M3 uses time. A tracker predicts where an instance should appear in the next frame, associates detections, and terminates a track when evidence is insufficient. Video-based oil-palm ripeness work shows why temporal input is attractive for deployment \cite{junior2023videopalm}. Temporal continuity is weaker when the camera moves around a tree and revisits a previously observed side. The report must state the maximum time gap and whether the camera path is ordered or arbitrary.

\subsection{Geometry and 3D association}

M4 associates observations through camera pose, depth, point clouds, or an SfM reconstruction. Gené-Molá et al. provide the clearest fruit example: a 2D instance mask is projected into a shared 3D representation so that repeated image detections can be consolidated \cite{genemola2020fruit3d}. Park et al. describe a newer real-time orchard approach that incrementally builds 3D crop instances and uses 3D generalized intersection-over-union for association \cite{park2026cropassociation}. The mechanism is directly aligned with duplicate rejection, but its field cost includes calibration, reconstruction stability, and sensor coverage.

Non-agricultural 3D detection studies provide transferable design components rather than fruit results. Camera and LiDAR fusion methods such as 3D-CVF and TransFusion show how multiple sensor views can be aligned at feature or decision level \cite{yoo20203dcvf,bai2022transfusion}. Visual SLAM systems show how a moving camera can estimate pose and maintain a scene map, while dynamic-scene variants explicitly mask moving objects before mapping \cite{bescos2018dynaslam}. These studies motivate a measurement question for agriculture: does a shared geometric frame reduce false duplicates after controlling for detector recall and sensor failure?

\subsection{Learned multi-view association}

M5 learns identity from linked observations. Training examples must state which detections belong to one physical object and which are different objects. Without this supervision, a model can learn a view-specific shortcut and still appear accurate on a random image split. SawitMVC's per-tree identity links create the type of supervision required for this experiment \cite{indriani2026sawitmvc}. The necessary split is by physical tree or plantation block, not by randomly distributing frames from one tree across train and test.

\section{Modalities and fusion choices}
\label{sec:fusion2}

RGB provides colour, texture, and fine appearance. Its weaknesses are shadows, specular highlights, camouflage, and scale changes. Depth provides a spatial cue, but ``depth'' is not one measurement. A calibrated sensor may provide metric range; stereo and SfM provide geometry relative to a reconstructed scene; a monocular network provides a learned or pseudo-depth ordering whose scale may not be metric. A manuscript must state which meaning is used.

Three fusion points are useful design categories. Early fusion concatenates RGB and depth at the input. It is simple but exposes the detector to missing or misregistered channels. Middle fusion extracts separate features and combines them at an intermediate layer, which permits modality-specific processing and quality gating. Late fusion combines detections, scores, or identity graphs after each modality has produced an estimate. RGB-depth object-detection studies demonstrate that the choice of fusion point is a measurable hypothesis rather than a universal rule \cite{ophoff2019multimodal}. The oil-palm RGB-depth CNN and visual-SLAM system illustrates a direct agricultural attempt to combine appearance, depth, and localization \cite{siong2021rgbdepthpalm}.

The design-space consequence is that sensor quality belongs in the evaluation, not only in the apparatus description. Each experiment should report the valid-depth ratio, RGB-depth registration error, missing-depth policy, and the fraction of detections supported by valid geometry. A fusion gain observed after dropping invalid frames is not comparable with a gain measured on all frames. The Goh dataset's reported registration error provides an example of the type of calibration evidence needed before interpreting a depth-based identity result \cite{goh2025palmdataset}.

\section{Synthesis: from detector score to inventory evidence}
\label{sec:synthesis2}

The evidence supports five design rules.

\begin{enumerate}[leftmargin=*,itemsep=2pt]
  \item Report per-view detection separately from cross-view identity. A detector score cannot establish duplicate-free counting.
  \item Define the unit of truth at tree level. For each physical tree, store the unique instance count, per-view appearances, class labels, and visibility state.
  \item Treat the acquisition path as an experimental factor. A fixed four-side census, a continuous video, and an arbitrary multi-camera setup impose different assumptions.
  \item Test identity mechanisms against their required assumptions. M1 needs regular duplication; M3 needs temporal overlap; M4 needs geometry; M5 needs identity-linked training data.
  \item Report class-wise count error after identity resolution. A correct class label on a duplicate is still a wrong inventory.
\end{enumerate}

The minimum evaluation bundle is therefore: per-instance precision and recall, class-wise average precision, identity precision and recall, duplicate rate, missed-appearance rate, absolute and relative count error per tree, class-confusion error per unique instance, sensor-validity statistics, and latency. Where tracks are used, IDF1 or HOTA can complement track precision and recall. Where 3D association is used, calibration error and association distance must be reported. These metrics should be paired by physical tree so that a difficult canopy is not assigned to one method and an easy canopy to another.

\subsection{Focused evidence matrix}

Appendix~\ref{app:matrix2} contains one row per focused study or review used to define the design space. The table includes direct agricultural evidence, fruit evidence, transferable mechanism evidence, and positioning reviews. A status field distinguishes a source with verified local full text from a source represented at the accessible article-record or abstract level. This distinction is part of the evidence, not an editorial footnote.

\section{Research gaps and testable agenda}
\label{sec:agenda2}

The first gap is an evaluation gap. Many studies report detection or ripeness classification, but a class-wise inventory requires identity links and a tree-level unique-count ground truth. SawitMVC demonstrates the data structure needed for this measurement, while the earlier fruit and FFB literature often reports image-level or frame-level results \cite{indriani2026sawitmvc,naftali2024palmcounter,koirala2019deepfruit}.

The second gap is a mechanism comparison gap. The literature contains examples of statistical correction, temporal video processing, RGB-depth fusion, and SfM, but these mechanisms are rarely compared under one acquisition protocol with the same detector and the same tree-level truth. The appropriate experiment is not ``model A versus model B'' alone. It is an M0-to-M4 or M5 ablation in which the detector, image split, class labels, and counting truth remain fixed.

The third gap is a sensor-assumption gap. A calibrated depth camera, a point cloud, LiDAR, and monocular pseudo-depth should not be grouped under one label. Each has a different scale, failure mode, and deployment cost. An outdoor FFB benchmark should report invalid depth, registration, range, illumination, and view baseline before assigning a performance gain to the modality.

The fourth gap is a class-resolution gap. Maturity is common, but the same identity record may need multiple attributes. A future benchmark should permit at least one class attribute beyond maturity or make clear why a single attribute is the intended use case. The evaluation should distinguish class confusion from identity duplication.

The resulting research agenda is a sequence of gates. Gate 1 verifies the per-view detector and the class labels. Gate 2 builds the unique-instance tree truth and measures the M0 baseline. Gate 3 tests M1 under the fixed view protocol. Gate 4 adds temporal or appearance association. Gate 5 adds calibrated geometry or registered depth. Gate 6 tests a learned association model with a tree-disjoint split. A mechanism advances only if it lowers tree-level count error without exceeding the latency and sensor-quality budget defined before the experiment.

\section{Limitations}
\label{sec:limitations2}

The search is reproducible from the captured raw exports and scripts, but the full-text stage is incomplete at the current cutoff. Scopus was available and Web of Science was not, so database coverage is not equivalent to a two-subscription search. OpenAlex metadata and abstracts are useful for discovery but do not replace article full text. The local audit found only six matches between the new master and the verified local PDF corpus. Consequently, the 20,035 full-text candidates are reported as a retrieval boundary, not as included studies.

The focused matrix also mixes direct agricultural studies, fruit studies, transferable mechanism papers, and reviews. This is deliberate for a design-space review, but the evidence classes are not interchangeable. A result from an autonomous-driving benchmark does not establish oil-palm performance. A review that discusses ripeness does not establish cross-view identity. A dataset descriptor establishes the availability of labels and sensors, not the superiority of a model.

The review does not estimate a pooled accuracy, rank all detectors, or claim that explicit deduplication is always beneficial. Its bounded claim is that cross-view identity is a necessary design variable for the stated inventory task and that the value of each identity mechanism should be evaluated under its acquisition assumptions.

\section{Conclusion}
\label{sec:conclusion2}

The target of agricultural multi-view perception is a unique, class-wise inventory, not a sum of independent detections. The revised protocol makes the search reproducible and exposes an important evidence boundary: a large candidate master has been screened automatically, but most candidate full texts are not yet locally available. The direct evidence includes oil-palm detector studies, multi-angle and RGB-depth datasets, fruit SfM, and orchard load estimation. Together they support a six-mechanism design space spanning intra-view counting, statistical correction, appearance, temporal tracking, geometry, and learned association.

The central contribution is a testable question rather than a fixed architecture: when does explicit identity resolution reduce tree-level duplicate and class-wise count error enough to justify its assumptions and cost? Answering that question requires tree-disjoint data, identity-linked annotations, sensor-quality reporting, and evaluation after aggregation. Oil palm is the principal case, but the design logic applies to other tree-fruit inventories whenever repeated observations must be converted into one count.

\appendix

\section{Focused evidence matrix}
\label{app:matrix2}

\small
\setlength{\LTleft}{0pt}
\setlength{\LTright}{0pt}
\setlength{\tabcolsep}{2pt}
\begin{longtable}{@{}>{\raggedright\arraybackslash}p{0.14\textwidth}>{\raggedright\arraybackslash}p{0.17\textwidth}>{\raggedright\arraybackslash}p{0.14\textwidth}>{\raggedright\arraybackslash}p{0.14\textwidth}>{\raggedright\arraybackslash}p{0.25\textwidth}>{\raggedright\arraybackslash}p{0.10\textwidth}@{}}
\caption{Row-level synthesis matrix. One row represents one study or review. Direct = agricultural or fruit evidence; transfer = non-agricultural mechanism evidence; review = positioning source.}\label{tab:matrix2}\\
\toprule
Study & Method and modality & Evaluation context & Identity or attribute & Finding or transfer boundary & Status\\
\midrule
\endfirsthead
\multicolumn{6}{c}{\tablename\ \thetable{} continued}\\
\toprule
Study & Method and modality & Evaluation context & Identity or attribute & Finding or transfer boundary & Status\\
\midrule
\endhead
\midrule
\multicolumn{6}{r}{Continued on next page}\\
\endfoot
\bottomrule
\endlastfoot
Redmon et al. \cite{redmon2016yolo} & One-stage RGB detector & PASCAL VOC; box detection & M0 & Establishes the speed-oriented per-image baseline; no repeated-view identity. & transfer\\
Ren et al. \cite{ren2017fasterrcnn} & Two-stage RGB detector with learned proposals & PASCAL VOC and COCO; box detection & M0 & Useful occlusion-accuracy comparator; no inventory identity. & transfer\\
He et al. \cite{he2017maskrcnn} & RGB instance segmentation & COCO; mask and box detection & M0 & Instance masks expose object separation before cross-view association. & transfer\\
Carion et al. \cite{carion2020detr} & Transformer set prediction in RGB & COCO; fixed prediction set & M0 & One-to-one matching removes duplicate predictions within an image, not duplicates across views. & transfer\\
Ge et al. \cite{ge2021yolox} & Anchor-free one-stage RGB detector & COCO; speed and AP & M0 & Supports a fast detector baseline; NMS-free or anchor-free does not guarantee cross-view uniqueness. & transfer\\
Wang et al. \cite{wang2024yolov10} & End-to-end RGB detector with one-to-one inference head & COCO; end-to-end detection & M0 & Removes an intra-image post-process; the physical identity problem remains. & transfer\\
Koirala et al. \cite{koirala2019deepfruit} & MangoYOLO and detector comparison, RGB video and images & Five orchards; fruit detection and load estimation & M1/M3 & High image detection still used a correction factor for orchard load; counting protocol affects interpretation. & direct\\
Gené-Molá et al. \cite{genemola2020fruit3d} & Mask R-CNN plus SfM point cloud & 11 Fuji trees and 1,455 apples & M4 & F1 increased from 0.816 in 2D to 0.881 for 3D detection and location; SfM was computationally costly. & full text\\
Fu et al. \cite{fu2020apple} & Faster R-CNN with RGB and depth features & Dense-foliage apple trees for robotic harvesting & M4 & Direct evidence for depth as an occlusion and localization cue; not an oil-palm benchmark. & direct\\
Kang and Chen \cite{kang2020strawberry} & Fruit detection, segmentation, and 3D visualization & Apple-orchard environment & M4 & Shows the value of 3D scene representation; identity and deployment cost require separate reporting. & direct\\
Siong et al. \cite{siong2021rgbdepthpalm} & RGB-depth CNN with heatmap localization and visual SLAM & Oil-palm FFB detection and localization & M4 & Direct palm example connecting multimodal detection to a mapping system; unique-count evaluation is not the primary reported unit. & metadata\\
Chen et al. \cite{chen2022ffbdet} & Computer-vision model comparison, RGB & Oil-palm FFB detection & M0 & Establishes a detector comparison context; no cross-view identity claim. & metadata\\
Lai et al. \cite{lai2022yolov4palm} & YOLOv4 for ripe FFB detection, RGB & Oil-palm ripeness detection & M0 & Direct maturity detector evidence; class assignment is not equivalent to duplicate-free inventory. & metadata\\
Junior and Suharjito \cite{junior2023videopalm} & Deep-learning model on video & Oil-palm ripeness video & M3 & Temporal input is relevant to tracking, but a video classifier is not automatically an identity graph. & metadata\\
Naftali et al. \cite{naftali2024palmcounter} & YOLO-family and state-of-the-art detector comparison & Plantation FFB dataset; detection and mobile-oriented counting & M0 & Establishes detector and efficiency baselines; repeated observations are not the main evaluation unit. & full text\\
Goh et al. \cite{goh2025palmdataset} & Registered RGB, depth maps, and point clouds & 400 RGB-depth pairs from four Johor locations & M4 & Reports 1.8 cm mean registration error at 3 m and 92.5\% expert agreement; supports sensor-quality reporting. & full text\\
Suharjito et al. \cite{suharjito2025palmdataset} & Smartphone video frames and object annotations & Plantation FFB dataset; five maturity classes & M0/M3 & Multiple angles and natural occlusion support inventory research; published split is not an identity benchmark by itself. & full text\\
Indriani et al. \cite{indriani2026sawitmvc} & Multi-angle smartphone RGB with identity-linked JSON & 3,992 images, 953 trees, four or eight views, 9,823 unique bunches & M1/M5 & Provides tree-level unique-bunch truth and cross-view links for direct deduplication experiments. & full text\\
Shiddiq et al. \cite{shiddiq2023palmcount} & Computer vision for FFB counting & Oil-palm counting task & M0/M1 & Direct count evidence; the identity mechanism and duplicate exposure must be checked in the full text. & metadata\\
Hamdani et al. \cite{hamdani2026palmcount} & Deep learning for natural-background FFB counting & Oil-palm yield estimation & M0/M1 & Recent search candidate for natural-background counting; full-text eligibility remains retrieval-dependent. & abstract\\
Park et al. \cite{park2026cropassociation} & Incremental 3D crop models with 3D GIoU association & Dense orchard counting & M4 & Directly targets duplicate reduction during acquisition; agricultural transfer is stronger than generic tracking but still not oil palm. & abstract\\
Yoo et al. \cite{yoo20203dcvf} & Camera-LiDAR feature fusion & Autonomous-driving 3D detection & M4/M5 & Provides a transferable fusion pattern; vehicle geometry is not fruit evidence. & transfer\\
Bai et al. \cite{bai2022transfusion} & Camera-LiDAR transformer fusion & Autonomous-driving 3D detection & M4/M5 & Shows learned cross-modal alignment; requires agricultural calibration and identity labels before transfer. & transfer\\
Bescos et al. \cite{bescos2018dynaslam} & RGB-D visual SLAM with dynamic-object handling & Indoor and dynamic-scene mapping & M4 & Supports separating mapping from moving objects; does not provide fruit-class or tree-level count evidence. & transfer\\
Ophoff et al. \cite{ophoff2019multimodal} & RGB-depth fusion comparison & Real-time object detection & Early/mid/late & Provides a framework for comparing fusion points; sensor-quality gating remains application-specific. & transfer\\
Lai et al. \cite{lai2023palmslr} & Systematic review of FFB ripeness methods & Oil-palm sensing literature & M0 & Positions ripeness sensors and methods; cross-view identity is not the central review unit. & review\\
Goh et al. \cite{goh2025palmreview} & Review of non-destructive FFB classification & Vision, LiDAR, spectroscopy, thermal, laser methods & M0/M4 & Positions modality trade-offs and real-time constraints; this review adds unique inventory identity. & review\\
Naghipour et al. \cite{naghipour2024fruitreview} & Review of AI fruit detection and classification & Agricultural fruit literature & M0 & Broad model and data survey; this review adds explicit cross-view identity as a design axis. & review\\
Sapkota et al. \cite{sapkota2024yoloagri} & Review of YOLO variants in agriculture & Agricultural applications and detector lineage & M0 & Positions YOLO in agriculture; this review is task-first rather than model-first. & review\\
Nordin and Misro \cite{nordin2026palmslr} & Systematic review of deep-learning palm maturity detection & Oil-palm maturity literature & M0 & Recent maturity-focused comparator; class-wise inventory is generalized beyond maturity here. & abstract\\
Zhai et al. \cite{zhai2025citrusyields} & Improved YOLOv8 with multi-period citrus data & Citrus yield estimation & M0/M1 & Broadens the agricultural comparison beyond palm; exact identity treatment requires full-text verification. & abstract\\
Ma et al. \cite{ma2026grapeyield} & YOLO-based grape detection and yield estimation & Grape yield framework & M0/M1 & Adds grape evidence to the search scope; it is not counted as cross-view proof without full-text review. & abstract\\
\end{longtable}

\section{Exact search strings}
\label{app:queries2}

The following blocks reproduce the final Scopus strings. All seven use \texttt{PUBYEAR > 2014 AND PUBYEAR < 2027} and the subject-area clause shown below. The line breaks are for typesetting only; the tokens are unchanged.

\subsection*{Q1: Multi-view fruit inventory}
\begin{verbatim}
TITLE-ABS-KEY(("fruit" OR "fresh fruit bunch" OR "FFB" OR "oil palm"
OR "apple" OR "citrus" OR "mango" OR "grape" OR "berry" OR "bunch"
OR "cluster" OR "crop")
AND ("count*" OR "yield estimation" OR "load estimation"
OR "fruit load" OR "inventory" OR "enumeration")
AND ("multi-view" OR "multiple view*" OR "multi-camera" OR "video"
OR "cross-view" OR "structure from motion" OR "SfM" OR "track*"
OR "3D reconstruction" OR "image sequence"))
AND PUBYEAR > 2014 AND PUBYEAR < 2027
AND (SUBJAREA(AGRI) OR SUBJAREA(COMP) OR SUBJAREA(ENGI)
OR SUBJAREA(MULT))
\end{verbatim}

\subsection*{Q2: Cross-view identity and duplicates}
\begin{verbatim}
TITLE-ABS-KEY(("multi-view" OR "cross-view" OR "multi-camera"
OR "multiple cameras" OR "overlapping view*" OR "multi-sensor")
AND ("re-identification" OR "data association" OR "identity"
OR "duplicate*" OR "deduplicat*" OR "double counting"
OR "correspondence" OR "instance matching" OR "track*"
OR "association")
AND ("object detection" OR "instance segmentation" OR "counting"
OR "instance" OR "pedestrian" OR "vehicle"))
AND PUBYEAR > 2014 AND PUBYEAR < 2027
AND (SUBJAREA(COMP) OR SUBJAREA(ENGI) OR SUBJAREA(AGRI)
OR SUBJAREA(MULT))
\end{verbatim}

\subsection*{Q3: Plant geometry and multimodal sensing}
\begin{verbatim}
TITLE-ABS-KEY(("RGB-D" OR "depth camera" OR "stereo vision" OR "LiDAR"
OR "point cloud" OR "monocular depth" OR "photogrammetry"
OR "structure from motion")
AND ("fruit" OR "orchard" OR "vineyard" OR "canopy" OR "plant"
OR "tree crop" OR "plantation" OR "oil palm")
AND ("detection" OR "segmentation" OR "localization" OR "counting"
OR "phenotyping" OR "harvesting"))
AND PUBYEAR > 2014 AND PUBYEAR < 2027
AND (SUBJAREA(AGRI) OR SUBJAREA(COMP) OR SUBJAREA(ENGI)
OR SUBJAREA(EART) OR SUBJAREA(ENVI) OR SUBJAREA(MULT))
\end{verbatim}

\subsection*{Q4: Per-instance attributes}
\begin{verbatim}
TITLE-ABS-KEY(("fruit" OR "produce" OR "crop" OR "bunch")
AND ("maturity" OR "ripeness" OR "grading" OR "grade"
OR "quality class*" OR "size class*" OR "defect" OR "disease"
OR "cultivar" OR "variety")
AND ("instance segmentation" OR "object detection" OR "per-instance"
OR "individual fruit" OR "bounding box" OR "per-object"))
AND PUBYEAR > 2014 AND PUBYEAR < 2027
AND (SUBJAREA(AGRI) OR SUBJAREA(COMP) OR SUBJAREA(ENGI)
OR SUBJAREA(MULT))
\end{verbatim}

\subsection*{Q5: Previous reviews}
\begin{verbatim}
TITLE(("review" OR "survey" OR "systematic review" OR "scoping review"))
AND TITLE-ABS-KEY(("fruit detection" OR "fruit counting"
OR "yield estimation" OR "oil palm" OR "orchard"
OR "crop monitoring" OR "agricultural robot*"
OR "precision agriculture")
AND ("deep learning" OR "computer vision" OR "object detection"
OR "machine vision"))
AND PUBYEAR > 2014 AND PUBYEAR < 2027
AND (SUBJAREA(AGRI) OR SUBJAREA(COMP) OR SUBJAREA(ENGI)
OR SUBJAREA(MULT))
\end{verbatim}

\subsection*{Q6: Oil-palm vision and machine learning}
\begin{verbatim}
TITLE-ABS-KEY(("oil palm" OR "elaeis guineensis"
OR "fresh fruit bunch" OR "FFB")
AND ("deep learning" OR "machine learning" OR "convolutional"
OR "neural network" OR "computer vision" OR "image processing"
OR "YOLO" OR "object detection" OR "remote sensing"))
AND PUBYEAR > 2014 AND PUBYEAR < 2027
AND (SUBJAREA(AGRI) OR SUBJAREA(COMP) OR SUBJAREA(ENGI)
OR SUBJAREA(EART) OR SUBJAREA(ENVI) OR SUBJAREA(MULT))
\end{verbatim}

\subsection*{Q7: Single-view fruit load}
\begin{verbatim}
TITLE-ABS-KEY(("fruit" OR "apple" OR "citrus" OR "mango" OR "grape"
OR "berry" OR "bunch" OR "oil palm" OR "orchard" OR "vineyard")
AND ("counting" OR "count" OR "yield estimation" OR "load estimation"
OR "fruit load" OR "crop load")
AND ("deep learning" OR "convolutional" OR "object detection"
OR "YOLO" OR "instance segmentation" OR "machine vision"))
AND PUBYEAR > 2014 AND PUBYEAR < 2027
AND (SUBJAREA(AGRI) OR SUBJAREA(COMP) OR SUBJAREA(ENGI)
OR SUBJAREA(MULT))
\end{verbatim}


\bibliographystyle{elsarticle-num}
\bibliography{references}
\end{document}
